%=============================================================================
% Chapter 2 Exercises - With hints from original book
%=============================================================================
\section*{Exercises}
\addcontentsline{toc}{section}{Exercises}

\begin{enumerate}
    \item Show that $(\Z/m\Z) \otimes_\Z (\Z/n\Z) = 0$ if $m, n$ are coprime.

    \item \label{exer2-chapter2}Let $A$ be a ring, $\mathfrak{a}$ an ideal, $M$ an $A$-module. Show that $(A/\mathfrak{a}) \otimes_A M$ is isomorphic to $M/\mathfrak{a}M$.\\
    \textit{Hint:} Tensor the exact sequence $0 \to \mathfrak{a} \to A \to A/\mathfrak{a} \to 0$ with $M$.

    \item \label{exer3-chapter2}Let $A$ be a local ring, $M$ and $N$ finitely generated $A$-modules. Prove that if $M \otimes N = 0$, then $M = 0$ or $N = 0$.\\
    \textit{Hint:} Let $\mathfrak{m}$ be the maximal ideal, $k = A/\mathfrak{m}$ the residue field. Let $M_k = k \otimes_A M \cong M/\mathfrak{m}M$ by Exercise \ref{exer2-chapter2}. By Nakayama's lemma, $M_k = 0 \implies M = 0$. But $M \otimes_A N = 0 \implies (M \otimes_A N)_k = 0 \implies M_k \otimes_k N_k = 0 \implies M_k = 0$ or $N_k = 0$, since $M_k, N_k$ are vector spaces over a field.

    \item \label{exer4-chapter2}Let $M_i$ ($i \in I$) be any family of $A$-modules, and let $M$ be their direct sum. Prove that $M$ is flat $\iff$ each $M_i$ is flat.

    \item \label{exer5-chapter2}Let $A[x]$ be the ring of polynomials in one indeterminate over a ring $A$. Prove that $A[x]$ is a flat $A$-algebra.\\
    \textit{Hint:} Use Exercise \ref{exer4-chapter2}.

    \item \label{exer6-chapter2}For any $A$-module, let $M[x]$ denote the set of all polynomials in $x$ with coefficients in $M$, that is to say expressions of the form $\sum_{i=0}^n m_i x^i$ ($m_i \in M$). Defining the product of an element of $A[x]$ and an element of $M[x]$ in the obvious way, show that $M[x]$ is an $A[x]$-module. 
    
    \quad\, Show that $M[x] \cong A[x] \otimes_A M$.

    \item \label{exer7-chapter2}Let $\mathfrak{p}$ be a prime ideal in $A$. Show that $\mathfrak{p}[x]$ is a prime ideal in $A[x]$. If $\mathfrak{m}$ is a maximal ideal in $A$, is $\mathfrak{m}[x]$ a maximal ideal in $A[x]$?

    \item \label{exer8-chapter2}\begin{enumerate}
        \item If $M$ and $N$ are flat $A$-modules, then so is $M \otimes_A N$.
        \item If $B$ is a flat $A$-algebra and $N$ is a flat $B$-module, then $N$ is flat as an $A$-module.
    \end{enumerate}

    \item \label{exer9-chapter2}Let $0 \to M' \to M \to M'' \to 0$ be an exact sequence of $A$-modules. If $M'$ and $M''$ are finitely generated, then so is $M$.

    \item \label{exer10-chapter2}Let $A$ be a ring $\neq 0$, $\mathfrak{a}$ an ideal contained in the Jacobson radical of $A$; let $M$ be an $A$-module and $N$ a finitely generated $A$-module, and let $u \colon M \to N$ be a homomorphism. If the induced homomorphism $M/\mathfrak{a}M \to N/\mathfrak{a}N$ is surjective, then $u$ is surjective.

    \item \label{exer11-chapter2}Let $A$ be a ring $\neq 0$. Show that $A^m \cong A^n \implies m = n$.\\
    \textit{Hint:} Let $\mathfrak{m}$ be a maximal ideal of $A$ and let $\phi \colon A^m \to A^n$ be an isomorphism. Then $1 \otimes \phi \colon (A/\mathfrak{m}) \otimes A^m \to (A/\mathfrak{m}) \otimes A^n$ is an isomorphism between vector spaces of dimensions $m$ and $n$ over the field $k = A/\mathfrak{m}$. Hence $m = n$. (Cf. Chapter 3, Exercise \ref{exer15-chapter3}.)\\
    \quad\, If $\phi \colon A^m \to A^n$ is surjective, then $m \geq n$.\\
    \quad\, If $\phi \colon A^m \to A^n$ is injective, is it always the case that $m \leq n$?

    \item \label{exer12-chapter2}Let $M$ be a finitely generated $A$-module and $\phi \colon M \to A^n$ a surjective homomorphism. Show that $\Ker(\phi)$ is finitely generated.\\
    \textit{Hint:} Let $e_1, \ldots, e_n$ be a basis of $A^n$ and choose $u_i \in M$ such that $\phi(u_i) = e_i$ ($1 \leq i \leq n$). Show that $M$ is the direct sum of $\Ker(\phi)$ and the submodule generated by $u_1, \ldots, u_n$.

    \item \label{exer13-chapter2}Let $f \colon A \to B$ be a ring homomorphism, and let $N$ be a $B$-module. Regarding $N$ as an $A$-module by restriction of scalars, form the $B$-module $N_B = B \otimes_A N$. Show that the homomorphism $g \colon N \to N_B$ which maps $y$ to $1 \otimes y$ is injective and that $g(N)$ is a direct summand of $N_B$.\\
    \textit{Hint:} Define $p \colon N_B \to N$ by $p(b \otimes y) = by$, and show that $N_B = \Image(g) \oplus \Ker(p)$.

    \emph{Direct limits.}

    \item \label{exer14-chapter2}A partially ordered set $I$ is said to be a \emph{directed set}\index{directed set} if for each pair $i, j$ in $I$ there exists $k \in I$ such that $i \leq k$ and $j \leq k$.\\
    Let $A$ be a ring, let $I$ be a directed set and let $(M_i)_{i \in I}$ be a family of $A$-modules indexed by $I$. For each pair $i, j$ in $I$ such that $i \leq j$, let $\mu_{ij} \colon M_i \to M_j$ be an $A$-homomorphism, and suppose that the following axioms are satisfied:
    \begin{enumerate}[(1)]
        \item $\mu_{ii}$ is the identity mapping of $M_i$, for all $i \in I$;
        \item $\mu_{ik} = \mu_{jk} \circ \mu_{ij}$ whenever $i \leq j \leq k$.
    \end{enumerate}
    Then the modules $M_i$ and homomorphisms $\mu_{ij}$ are said to form a \emph{direct system}\index{direct system} $\mathbf{M} = (M_i, \mu_{ij})$ over the directed set $I$.

    \quad\, We shall construct an $A$-module $M$ called the \emph{direct limit}\index{direct limit} of the direct system $\mathbf{M}$. Let $C$ be the direct sum of the $M_i$, and identify each module $M_i$ with its canonical image in $C$. Let $D$ be the submodule of $C$ generated by all elements of the form $x_i - \mu_{ij}(x_i)$ where $i \leq j$ and $x_i \in M_i$. Let $M = C/D$, let $\mu \colon C \to M$ be the projection and let $\mu_i$ be the restriction of $\mu$ to $M_i$.

    \quad\, The module $M$, or more correctly the pair consisting of $M$ and the family of homomorphisms $\mu_i \colon M_i \to M$, is called the direct limit of the direct system $\mathbf{M}$, and is written $\varinjlim M_i$. From the construction it is clear that $\mu_i = \mu_j \circ \mu_{ij}$ whenever $i \leq j$.

    \item \label{exer15-chapter2}In the situation of Exercise 14, show that every element of $M$ can be written in the form $\mu_i(x_i)$ for some $i \in I$ and some $x_i \in M_i$. Show also that if $\mu_i(x_i) = 0$ then there exists $j \geq i$ such that $\mu_{ij}(x_i) = 0$ in $M_j$.

    \item \label{exer16-chapter2}Show that the direct limit is characterized (up to isomorphism) by the following property. Let $N$ be an $A$-module and for each $i \in I$ let $\alpha_i \colon M_i \to N$ be an $A$-module homomorphism such that $\alpha_i = \alpha_j \circ \mu_{ij}$ whenever $i \leq j$. Then there exists a unique homomorphism $\alpha \colon M \to N$ such that $\alpha_i = \alpha \circ \mu_i$ for all $i \in I$.

    \item \label{exer17-chapter2}Let $(M_i)_{i \in I}$ be a family of submodules of an $A$-module, such that for each pair of indices $i, j$ in $I$ there exists $k \in I$ such that $M_i + M_j \subseteq M_k$. Define $i \leq j$ to mean $M_i \subseteq M_j$ and let $\mu_{ij} \colon M_i \to M_j$ be the embedding of $M_i$ in $M_j$. Show that
    \[
    \varinjlim M_i = \sum M_i = \bigcup M_i.
    \]
    In particular, any $A$-module is the direct limit of its finitely generated submodules.

    \item \label{exer18-chapter2}Let $\mathbf{M} = (M_i, \mu_{ij})$, $\mathbf{N} = (N_i, \nu_{ij})$ be direct systems of $A$-modules over the same directed set. Let $M, N$ be the direct limits and $\mu_i \colon M_i \to M$, $\nu_i \colon N_i \to N$ the associated homomorphisms. 
    
    A \emph{homomorphism} $\Phi \colon \mathbf{M} \to \mathbf{N}$ is by definition a family of $A$-module homomorphisms $\phi_i \colon M_i \to N_i$ such that $\phi_j \circ \mu_{ij} = \nu_{ij} \circ \phi_i$ whenever $i \leq j$. Show that $\Phi$ defines a unique homomorphism $\phi = \varinjlim \phi_i \colon M \to N$ such that $\phi \circ \mu_i = \nu_i \circ \phi_i$ for all $i \in I$.

    \item \label{exer19-chapter2}A sequence of direct systems and homomorphisms $\mathbf{M} \to \mathbf{N} \to \mathbf{P}$ is exact if the corresponding sequence of modules and module homomorphisms is exact for each $i \in I$. Show that the sequence $\varinjlim \mathbf{M} \to \varinjlim \mathbf{N} \to \varinjlim \mathbf{P}$ of direct limits is then exact.\\
    \textit{Hint:} Use Exercise \ref{exer15-chapter2}.

    \emph{Tensor products commute with direct limits.} 

    \item \label{exer20-chapter2}Keeping the same notation as in Exercise 14, let $N$ be any $A$-module. Then $(M_i \otimes N, \mu_{ij} \otimes 1)$ is a direct system; let $P = \varinjlim (M_i \otimes N)$ be its direct limit.
    
    For each $i \in I$ we have a homomorphism $\mu_i \otimes 1 \colon M_i \otimes N \to M \otimes N$, hence by Exercise 16 a homomorphism $\psi \colon P \to M \otimes N$. Show that $\psi$ is an isomorphism, so that
    \[
    \varinjlim (M_i \otimes N) \cong (\varinjlim M_i) \otimes N.
    \]
    \textit{Hint:} For each $i \in I$, let $g_i \colon M_i \times N \to M_i \otimes N$ be the canonical bilinear mapping. Passing to the limit we obtain a mapping $g \colon M \times N \to P$. Show that $g$ is $A$-bilinear and hence define a homomorphism $\phi \colon M \otimes N \to P$. Verify that $\phi \circ \psi$ and $\psi \circ \phi$ are identity mappings.

    \item \label{exer21-chapter2}Let $(A_i)_{i \in I}$ be a family of rings indexed by a directed set $I$, and for each pair $i \leq j$ in $I$ let $\alpha_{ij} \colon A_i \to A_j$ be a ring homomorphism, satisfying conditions (1) and (2) of Exercise \ref{exer14-chapter2}. Regarding each $A_i$ as a $\Z$-module we can then form the direct limit $A = \varinjlim A_i$. Show that $A$ inherits a ring structure from the $A_i$ so that the mappings $A_i \to A$ are ring homomorphisms. The ring $A$ is the \emph{direct limit} of the system $(A_i, \alpha_{ij})$.
    
    \quad\, If $A = 0$ prove that $A_i = 0$ for some $i \in I$.

    \textit{Hint:} Remember that all rings have identity elements!

    \item \label{exer22-chapter2}Let $(A_i, \alpha_{ij})$ be a direct system of rings and let $\mathfrak{N}_i$ be the nilradical of $A_i$. Show that $\varinjlim \mathfrak{N}_i$ is the nilradical of $\varinjlim A_i$.
    
    If each $A_i$ is an integral domain, then $\varinjlim A_i$ is an integral domain.

    \item \label{exer23-chapter2}Let $(B_\lambda)_{\lambda \in \Lambda}$ be a family of $A$-algebras. For each finite subset of $\Lambda$ let $B_J$ denote the tensor product (over $A$) of the $B_\lambda$ for $\lambda \in J$. If $J'$ is another finite subset of $\Lambda$ and $J \subseteq J'$, there is a canonical $A$-algebra homomorphism $B_J \to B_{J'}$. Let $B$ denote the direct limit of the rings $B_J$ as $J$ runs through all finite subsets of $\Lambda$. The ring $B$ has a natural $A$-algebra structure for which the homomorphisms $B_J \to B$ are $A$-algebra homomorphisms. The $A$-algebra $B$ is the \emph{tensor product}\index{tensor product} of the family $(B_\lambda)_{\lambda \in \Lambda}$.

    \emph{Flatness and Tor.} 

    In these Exercises it will be assumed that the reader is familiar with the definition and basic properties of the Tor functor.

    \item \label{exer24-chapter2}If $M$ is an $A$-module, the following are equivalent:
    \begin{enumerate}[i)]
        \item $M$ is flat;
        \item $\Tor_n^A(M, N) = 0$ for all $n > 0$ and all $A$-modules $N$;
        \item $\Tor_1^A(M, N) = 0$ for all $A$-modules $N$.
    \end{enumerate}
    \textit{Hint:} To show that (i) $\implies$ (ii), take a free resolution of $N$ and tensor it with $M$. Since $M$ is flat, the resulting sequence is exact and therefore its homology groups, which are the $\Tor_n^A(M, N)$, are zero for $n > 0$. To show that (iii) $\implies$ (i), let $0 \to N' \to N \to N'' \to 0$ be an exact sequence. Then, from the Tor exact sequence,
    \[
    \Tor_1^A(M, N'') \to M \otimes N' \to M \otimes N \to M \otimes N'' \to 0
    \]
    is exact. Since $\Tor_1^A(M, N'') = 0$ it follows that $M$ is flat.

    \item \label{exer25-chapter2}Let $0 \to N' \to N \to N'' \to 0$ be an exact sequence, with $N''$ flat. Then $N'$ is flat $\iff$ $N$ is flat.

    \textit{Hint:} Use Exercise \ref{exer24-chapter2} and the Tor exact sequence.

    \item \label{exer26-chapter2}Let $N$ be an $A$-module. Then $N$ is flat $\iff$ $\Tor_1^A(A/\mathfrak{a}, N) = 0$ for all finitely generated ideals $\mathfrak{a}$ in $A$.\\
    \textit{Hint:} Show first that $N$ is flat if $\Tor_1^A(M, N) = 0$ for all finitely generated $A$-modules $M$, by using \eqref{prop:flat-equivalent}. If $M$ is finitely generated, let $x_1, \ldots, x_n$ be a set of generators of $M$, and let $M_i$ be the submodule generated by $x_1, \ldots, x_i$. By considering the successive quotients $M_i/M_{i-1}$ and using Exercise \ref{exer25-chapter2}, deduce that $N$ is flat if $\Tor_1^A(M, N) = 0$ for all \emph{cyclic}\index{cyclic} $A$-modules $M$, i.e., all $M$ generated by a single element, and therefore of the form $A/\mathfrak{a}$ for some ideal $\mathfrak{a}$. Finally use \eqref{prop:flat-equivalent} again to reduce to the case where $\mathfrak{a}$ is a finitely generated ideal.

    \item \label{exer27-chapter2}A ring $A$ is \emph{absolutely flat}\index{absolutely flat ring} if every $A$-module is flat. Prove that the following are equivalent:
    \begin{enumerate}[i)]
        \item $A$ is absolutely flat.
        \item Every principal ideal is idempotent.
        \item Every finitely generated ideal is a direct summand of $A$.
    \end{enumerate}
    \textit{Hint:} (i) $\implies$ (ii). Let $x \in A$. Then $A/(x)$ is a flat $A$-module, hence in the diagram
    \[\begin{tikzcd}
        {(x)\otimes A} & {(x)\otimes A/(x)} \\
        A & {A/(x)}
        \arrow["\beta", from=1-1, to=1-2]
        \arrow[from=1-1, to=2-1]
        \arrow["\alpha", from=1-2, to=2-2]
        \arrow[from=2-1, to=2-2]
    \end{tikzcd}\]
    the mapping $\alpha$ is injective. Hence $\Image(\beta) = 0$, hence $(x) = (x^2)$.\\
    (ii) $\implies$ (iii). Let $x \in A$. Then $x = ax^2$ for some $a \in A$, hence $e = ax$ is idempotent and we have $(e) = (x)$. Now if $e, f$ are idempotents, then $(e, f) = (e + f - ef)$. Hence every finitely generated ideal is principal, and generated by an idempotent $e$, hence is a direct summand because $A = (e) \oplus (1-e)$.\\
    (iii) $\implies$ (i). Use the criterion of Exercise \ref{exer26-chapter2}.

    \item \label{exer28-chapter2}A Boolean ring is absolutely flat. The ring of Chapter 1, Exercise \ref{exer7-chapter1} is absolutely flat. Every homomorphic image of an absolutely flat ring is absolutely flat. If a local ring is absolutely flat, then it is a field.
    
    \quad\, If $A$ is absolutely flat, every non-unit in $A$ is a zero-divisor.
\end{enumerate}
